\section{Introduction}
\label{sec:intro}

The widespread deployment of large foundation models on resource-constrained edge devices (e.g., mobile phones, robotic platforms, and IoT hardware) has revolutionized distributed artificial intelligence. However, edge-deployed models are frequently subjected to highly dynamic, non-stationary task streams where requests from completely different domains (such as text, image, speech, and medical diagnostics) arrive in real-time. Maintaining a distinct, full-sized large neural network for each individual task is computationally and financially prohibitive. To bypass this, parameter-efficient fine-tuning (PEFT) methods, particularly Low-Rank Adaptation (LoRA) \cite{hu21lora}, are employed to specialize a single pre-trained backbone model to various tasks. At test-time, these specialized adapter weights are combined on-the-fly to handle diverse input streams, a process known as \textit{dynamic model merging} or test-time ensembling \cite{liang23sable, zhang24spszca}.

Despite the promise of dynamic model merging, current state-of-the-art serving systems encounter a severe stability-accuracy bottleneck. Standard stateless ensembling methods (e.g., SABLE \cite{liang23sable}) compute routing similarities independently at each layer. Because deep neural network representations can fluctuate significantly from layer to layer, these stateless approaches suffer from catastrophic layer-to-layer ensembling weight jitter. This weight jitter means that the model's blended parameters oscillate rapidly through weight space across sequential depth, disrupting the smooth propagation of activations and causing representational incoherence. To address this, state-dependent merging systems like ChemMerge have proposed using first-order non-equilibrium chemical reaction kinetics to smooth routing weight updates. However, first-order chemical ODE systems lack physical inertia and momentum; under competitive reaction rates (e.g., low temperature regimes), they are highly prone to volatile concentration oscillations and overreactions, amplifying routing weight jitter rather than resolving it.

In this work, we propose a novel alternative to both stateless and chemical ensembling paradigms. Drawing inspiration from classical mechanics and orbital dynamics, we present \textbf{GraviMerge}, a mathematically rigorous, physics-informed model merging framework that models deep representation routing as a multi-body physical system. We argue that deep activation trajectories can be modeled and stabilized through the elegant, universal laws of physical dynamics. As illustrated in Figure~\ref{fig:layer_trajectory}, GraviMerge represents the intermediate activation probe flowing through network layers as a virtual stateful spacecraft. The pre-trained expert centroids are modeled as stationary, high-mass attractors exerting softened gravitational forces. Crucially, by modeling representation flow with virtual physical mass, velocity, acceleration, and viscous medium drag, we introduce true second-order inertia directly into the routing trajectory. 

Furthermore, to maintain absolute mathematical and geometric consistency on the curved representation manifold, we incorporate exact spherical geodesic updates (exponential mapping) and parallel transport of the velocity state vector. This prevents numerical energy accumulation and drift, providing an exceptionally robust and mathematically sound system. Because the spacecraft probe has physical momentum and travels along the manifold's natural geodesics, any high-frequency representational noise is naturally filtered and smoothed, resulting in highly stable ensembling weight trajectories across deep layers.

We make the following core contributions:
\begin{itemize}
    \item \textbf{A Physics-Informed Routing Paradigm: } We introduce GraviMerge, a dynamic model merging framework that models activation trajectories using multi-body gravitational physics on spherical manifolds. This offers an interpretable and robust formulation of representational stability.
    \item \textbf{Novel Physical and Geometric Mechanisms: } We formulate three novel routing mechanisms: (1) \textit{Arrhenius Mass Activation (AMA)} to dynamically set expert masses via early zero-shot alignment; (2) \textit{Geodesic Trajectory Integration (GTI)} which integrates gravitational force, velocity, and viscous drag to update the probe's coordinate trajectory via exact spherical geodesic steps and parallel transport; and (3) \textit{Gravitational Influence Blending (GIB)} to translate localized gravitational pull into continuous ensembling weights.
    \item \textbf{Rigorous Stability and Robustness: } We perform exhaustive evaluations across $10$ independent seeds on a Projected Digit Representation Space (RDS) Proxy benchmark. GraviMerge completely resolves the accuracy-stability dilemma, achieving the highest joint serving accuracy (\textbf{88.69\%}) while dramatically reducing layer-to-layer ensembling weight jitter by \textbf{6.01$\times$} compared to ChemMerge, \textbf{5.47$\times$} compared to first-order EMA, and \textbf{2.40$\times$} compared to SABLE, completely outclassing weight-space second-order momentum filters (which increase jitter by \textbf{14.54$\times$}).
    \item \textbf{Resilience to Dynamic Workloads: } We demonstrate that GraviMerge is highly resilient to dynamic edge-serving workloads, maintaining optimal performance across homogeneous, heterogeneous, and vectorized ($B=1$) serving environments, enabling robust and computationally efficient real-time deployment.
\end{itemize}
